% claim-ref: eisenia-fetida#biology
% claim-ref: eisenia-fetida#cocoon-timing
% claim-ref: eisenia-fetida#colony-protocol
% claim-ref: eisenia-fetida#harvest-surplus
\CompanionHeader{REPRODUCTION}{Cocoons and colony care}{Red wiggler / Eisenia fetida / planning only}{Feasible in stable bins; success must be observed, not inferred from elapsed time.}
\Context{Simultaneous hermaphrodites usually exchange sperm; the clitellum forms cocoons. Laboratory self-fertilization and fertile E. andrei hybrids are documented: no absolute impossibility claim. Keep identity uncertainty in the record. [EF-N2; EF-T2]}
\Step{1}{Prepare parents and capacity}{Traceable stock, prepared bedding and a record of mature clitellate adults first. Stabilize consumption, moisture and temperature; arrange capacity for surplus before encouraging expansion. [EF-N1; EF-N2; proposed]}
\Step{2}{Keep gentle care}{Leave cocoons in moist, aerated bedding; avoid frequent sorting. Bedding and small decomposing food portions remain accessible to the colony. Do not add excess feed to force juvenile growth. [EF-N1; EF-O1; proposed]}
\Step{3}{Observe distinct milestones}{Feeding adults, cocoons, live juveniles, growing juveniles, then mature worms. Yellowish cocoons darken. NC State examples: about three hatchlings per cocoon; 30--75 days to hatch, then another 53--76 days to maturity. Conditional ranges, not local deadlines; a short wait does not prove cocoon death. [EF-N2]}
\Step{4}{Pause when conditions deteriorate}{Accumulating food, rotten odor, unsuitable temperature or repeated escapes stop expansion while conditions are corrected. Do not compensate with extra feed. Before splitting, prepare both bins and retain bedding, adults and juveniles. No promised population doubling. [EF-N1; EF-U1; proposed]}
\Step{5}{Harvest gently}{When material is dark/crumbly and bedding needs renewal, offer fresh bedding/food on one side; allow migration, then remove finished material from the other. Return remaining worms and some old material to fresh bedding. UC gives about 2.5--6 months, varying with conditions. [EF-U1]}
\Step{6}{Contain surplus and verify transfer}{No intentional release of worms, cocoons or unsorted culture into natural areas, drains or waterways. Keep suitable extra capacity or arrange an identified contained transfer; otherwise ask local Extension for humane disposal and waste advice. Verify origin/destination requirements first. Minnesota prevention guidance is a precautionary analogy, not California law or an E. fetida invasive designation. [EF-M1; proposed]}
\Note{Observation record}{Cocoon / juvenile first seen; approximate counts; growth: \hrulefill\par
Feed / bedding temperature / moisture / capacity / next review: \hrulefill}
\Note{Shared handwritten log}{Choose condition-based moisture care, not scheduled pouring. Use A/M for feed or misting and Obs for cocoons/juveniles; leave unperformed actions blank.}
\Sources{Clickable evidence: \href{https://content.ces.ncsu.edu/worms-can-recycle-your-garbage}{EF-N1}; \href{https://content.ces.ncsu.edu/raising-earthworms-successfully}{EF-N2}; \href{https://extension.oregonstate.edu/catalog/em-9034-composting-worms}{EF-O1}; \href{https://ucanr.edu/node/161011/printable/print}{EF-U1}; \href{https://extension.umn.edu/natural-resources/forestry-and-wildlife/invasive-species/jumping-worms}{EF-M1}; \href{https://journals.plos.org/plosone/article?id=10.1371/journal.pone.0191711}{EF-T2}. Claims, source scope, separate dates and contradictions: adjacent YAML. Planning only; habitat unmeasured.}{eisenia-fetida / reproduction}
